% Folder 93, batch 05 — archivists' pages 81–100.
% Pass: Opus 5.5 (claude-opus-5-5), 2026-10-10 — first pass, unchecked against the pages by a human.
%
% Page 84 is a blank sheet and is skipped.
\documentclass[11pt,a4paper]{article}
\input{../preamble/grothendieck}

\folder{93}
\batch{5}
\pages{81}{100}
\foldertitle{Connexions de Gauss-Manin et équations de Picard-Fuchs. Opération de Cartier : copies de tapuscrit annoté (s.d.), tiré à part (1981), notes manuscrites (s.d.).}
\dating{[1972]-1981}
\watermark{Édition de démonstration}

\begin{document}

\page{81}
\note{feuille de calculs, en deux parties séparées par un trait horizontal ; la première est un schéma de flèches entre faisceaux quotients de $\mathcal{E}_X$, la seconde une liste d'« écarts » suivie d'une seconde version du même schéma. Les flèches courbes verticales du schéma sont décrites plutôt que redessinées.}

« écart additif »
\[
A_1(f,g) = dg/df = \frac{g'}{f'}, \qquad
A_n(f,g) = \Bigl(\frac{dg}{df}\Bigr)^n = \Bigl(\frac{g'}{f'}\Bigr)^n
\]

« écart affine »
\[
A_{\mathrm{af}}(f,g) = D_{dg} - D_{df}
= \frac{d\,\frac{dg}{df}}{dg/df}
= \frac{d\,\frac{g'}{f'}}{g'/f'}
= \frac{dg'}{g'} - \frac{df'}{f'}
= \Bigl(\frac{g''}{g'} - \frac{f''}{f'}\Bigr)dt,
\qquad c_{\mathrm{af}}(f) = D_{df}
\]

\begin{tikzcd}[column sep=small, row sep=small, nodes={font=\scriptsize}]
\mathcal{E}_X \arrow[r] \arrow[d, "\exp"'] & \mathcal{E}_X/\mathbb{G}_a(\mathbb{C}) \overset{\sim}{\rightarrow} \underline{\omega}^{1*}_X \arrow[r, "\varphi\mapsto\varphi^n"] & \mathcal{E}_X/(\mu_n\cdot\mathbb{G}_a)(\mathbb{C}) \overset{\sim}{\rightarrow} \underline{\omega}^{n}_X \arrow[r, "{\psi\mapsto -\frac{1}{n}\frac{d\psi}{\psi}}"] & \mathcal{E}_X/(\mathbb{G}_m\cdot\mathbb{G}_a)(\mathbb{C}) \overset{\sim}{\rightarrow} \mathrm{Conn}(\underline{\omega}^1_X) \arrow[d, "D\mapsto c_D"] \\
\mathcal{E}^*_X \arrow[r] & \mathcal{E}^*_X/\mathbb{G}_m(\mathbb{C}) \overset{\sim}{\rightarrow} \underline{\omega}^{1*}_X \arrow[r, "\varphi\mapsto\varphi^2"] & \mathcal{E}^*_X/(\mathbb{Z}/2\mathbb{Z}\cdot\mathbb{G}_m)(\mathbb{C}) \simeq \underline{\omega}^{2*}_X & \mathcal{E}_X/\mathrm{GP}(1,\mathbb{C}) \overset{\sim}{\rightarrow} \mathrm{Connproj}(X)
\end{tikzcd}

\note{sous chaque quotient il écrit le groupe qui opère : $\underline{\omega}^{1*}_X$ et $\underline{\omega}^{n}_X$ « torseur sous $\underline{O}^*_X$ » ($n\neq 0$), $\mathrm{Conn}(\underline{\omega}^1_X)$ « torseur sous $\underline{\omega}^1_X$ », $\mathrm{Connproj}(X)$ « ($\underline{\omega}^2_X$ torseur sous) », $\underline{\omega}^{1*}_X$ (ligne du bas) « torseur sous $\underline{O}^*_X$ », $\underline{\omega}^{2*}_X$ « torseur sous $\underline{O}^{*}_X$ ». Les applications de classe sont notées sur les flèches : $c_a(f) = df$, $c_{a,n}(f) = (df)^n$, $c_m(f) = \frac{df}{f}$, $c_{\mathrm{pr}}(f) = c_{D_{df}}$. Deux flèches courbes, marquées « 2 » et « id », descendent de $\underline{\omega}^{1*}_X$ (ligne du haut) vers la ligne du bas. Le groupe $\mathbb{Z}/2\mathbb{Z}$ est une lecture \uncertain{probable} d'un tracé qui ressemble à « $\mathbb{Z}.n\mathbb{Z}$ » ; la lettre de $\mathcal{E}_X/\mathrm{GP}(1,\mathbb{C})$ est d'un tracé différent des autres $\mathcal{E}$.}

Dans un cadre :
\[
c_{D+\omega_1} = c_D + 2D\omega_1 + \omega_1^2 .
\]

\[
\exp(\pm f + b) = \exp b\,(\exp f)^{\pm 1}
\]
\note{le signe devant $f$ est une lecture \uncertain{probable}.}

\[
\begin{aligned}
\text{Écad}(f,g) &= \frac{dg}{df} = \frac{g'}{f'} \in \Gamma\,\underline{O}^*_X\\
\text{Écaf}_n(f,g) &= \Bigl(\frac{dg}{df}\Bigr)^n = \Bigl(\frac{g'}{f'}\Bigr)^n \in \Gamma(\underline{O}^*_X)\\
\text{Écaf}(f,g) &= \frac{d(dg/df)}{dg/df} = \frac{d(g'/f')}{g'/f'} = \frac{dg'}{g'} - \frac{df'}{f'}\\
&= \Bigl(\frac{g''}{g'} - \frac{f''}{f'}\Bigr)dt \in \Gamma(\underline{\omega}^1_X)\\
\text{Écmu}(f,g) &= \frac{dg/g}{df/f} = \frac{g'/g}{f'/f} \in \Gamma(\underline{O}^*_X)\\
\text{Écmu}_2(f,g) &= \Bigl(\frac{dg/g}{df/f}\Bigr)^2 = \Bigl(\frac{g'/g}{f'/f}\Bigr)^2 \in \Gamma(\underline{O}^*_X)\\
\text{Écpr}(f,g) &= -2\,\frac{d^3g/df^3}{dg/df} + 3\Bigl(\frac{d^2g/df^2}{dg/df}\Bigr)^2\\
&= \Bigl[\Bigl(-2\frac{g'''}{g'} + 3\Bigl(\frac{g''}{g'}\Bigr)^2\Bigr) - \Bigl(-2\frac{f'''}{f'} + 3\Bigl(\frac{f''}{f'}\Bigr)^2\Bigr)\Bigr]dt^2\\
&\in \Gamma(\underline{\omega}^2_X)
\end{aligned}
\]

Seconde version du schéma :

\begin{tikzcd}[column sep=small, row sep=small, nodes={font=\scriptsize}]
\mathcal{E}_X \arrow[r, "d"] \arrow[d, "\exp"'] & {[\mathcal{E}_X/\mathbb{G}_a(\mathbb{C}) \overset{\sim}{\rightarrow} \underline{\omega}^{1*}_X]} \arrow[r, "\alpha_n"] & {[\mathcal{E}_X/(\mu_n\cdot\mathbb{G}_a)(\mathbb{C}) \overset{\sim}{\rightarrow} \underline{\omega}^{n*}_X]} \arrow[r, "\beta_n"] & {[\mathcal{E}_X/(\mathbb{G}_m\cdot\mathbb{G}_a)(\mathbb{C}) \overset{\sim}{\rightarrow} \mathrm{Conn}(\underline{\omega}^1_X)]} \arrow[d, "\gamma"] \\
\mathcal{E}^*_X \arrow[r] & {[\mathcal{E}^*_X/\mathbb{G}_m(\mathbb{C}) \simeq \underline{\omega}^{1*}_X]} \arrow[r, "\alpha'_2"] & {[\mathcal{E}^*_X/(\mathbb{Z}/2\mathbb{Z}\cdot\mathbb{G}_m)(\mathbb{C}) \simeq \underline{\omega}^{2*}_X]} & {[\mathcal{O}_X/\mathrm{GP}(1)(\mathbb{C}) \simeq \mathrm{Conn}(X)]}
\end{tikzcd}

\note{sous la ligne du haut, les groupes opérants et leurs flèches : $\underline{O}^*_X \xrightarrow{\varphi\mapsto\varphi^n} \underline{O}^*_X \xrightarrow{\psi\mapsto -\frac{1}{n}\frac{d\psi}{\psi}} \underline{\omega}^1_X$ ; sous la ligne du bas, $\underline{O}^*_X \xrightarrow{\varphi\mapsto\varphi^2} \underline{O}^*_X \to \underline{\omega}^2_X$. Des flèches verticales « id », « 2 », « id ($n=2$) » relient les deux étages, et une longue flèche courbe $\delta$ part de la ligne du bas vers $\underline{\omega}^2_X$. Le terme $[\mathcal{O}_X/\mathrm{GP}(1)(\mathbb{C}) \simeq \mathrm{Conn}(X)]$ est une lecture \uncertain{douteuse}.}

\[
\begin{aligned}
\alpha_n(\omega_1) &= \omega_1^n\\
\beta_n(\underset{\substack{\|\\ f\,dt^n}}{\omega_n}) &= \beta_1(\omega_n^{1/n}) = \beta_1(f^{1/n}\,dt) = D_{dt} + \frac{1}{n}\frac{df}{f} = D_{dt} + \frac{1}{n}\frac{f'}{f}\,dt\\
\gamma(D_{dt} + \underset{\substack{\|\\ f\,dt}}{\omega_1}) &= \underset{\substack{\|\\ c_{\mathrm{pr}}(t)}}{\gamma(D_{dt})} + 2D_{dt}(\omega_1) + \omega_1^2\\
&= c_{\mathrm{pr}}(t) + 2\frac{df}{dt}\,dt + f^2\,dt^2 = c_{\mathrm{pr}}(t) + (2f' + f^2)\,dt^2\\
\gamma\beta_n(\omega_n) &= c_{\mathrm{pr}}(t) + \Bigl(2\Bigl(\frac{1}{n}\frac{f'}{f}\Bigr)' + \Bigl(\frac{1}{n}\frac{f'}{f}\Bigr)^2\Bigr)dt^2\\
&= c_{\mathrm{pr}}(t) + \frac{1}{n^2}\Bigl(-2n\frac{f''}{f}\ \ldots\ \Bigl(\frac{f'}{f}\Bigr)^2\Bigr)dt^2\\
&= c_{\mathrm{pr}}(t) + \Bigl(-2\frac{f''}{f} + 3\Bigl(\frac{f'}{f}\Bigr)^2\Bigr)dt^2\\
\gamma\beta_1\,d(f) &= c_{\mathrm{pr}}(t) + \Bigl(-2\frac{f'''}{f'} + 3\Bigl(\frac{f''}{f'}\Bigr)^2\Bigr)dt^2
\end{aligned}
\]
\note{la deuxième ligne de $\gamma\beta_n(\omega_n)$ est surchargée (un coefficient du type « $2n+1$ » est écrit par-dessus) et n'est lue qu'en partie : \ill{}. Une ligne « $\gamma\beta_1(\omega_1)$ » est commencée puis laissée, et un terme \struck{\ill{}} est biffé dans la dernière ligne.}

Encadré à droite :
\[
\begin{aligned}
\delta(\omega_2) &= \gamma\beta_2(\omega_2) + \omega_2\\
&= c_{\mathrm{pr}}(t) + \Bigl[\ldots\frac{f''}{f} + \frac{5}{4}\Bigl(\frac{f'}{f}\Bigr)^2\Bigr]dt^2\\
\delta\alpha'_2(\omega_1) &= \gamma\beta_1(\omega_1) + \omega_1^2\\
&= c_{\mathrm{pr}}(t) + \Bigl[f^2 - 2\frac{f''}{f} + 3\Bigl(\frac{f'}{f}\Bigr)^2\Bigr]dt^2
\end{aligned}
\]
\note{encadré très surchargé : le premier coefficient de $\delta(\omega_2)$ est \ill{}, le $\frac{5}{4}$ est \uncertain{douteux}, et la dernière ligne n'est lue que \uncertain{approximativement}.}

\page{82}
\note{feuille isolée, sur les invariants de Postnikov ; elle ne se rattache pas visiblement aux calculs qui l'entourent.}

$\uncertain{X} = \Omega Y$

\[
\begin{aligned}
&k_1 \in H^3(\pi_1(X), \pi_2(X)) \overset{?}{=} 0\\
&\qquad\uparrow \qquad\qquad \uparrow\\
\exists\ &k_1' \in H^4(K(\pi_1(X),2), \pi_2(X))\\
&\qquad\qquad \uparrow\\
&k_1'' \quad H^5(K(\pi_1(X),3), \pi_2(X))\\
&\qquad\qquad \simeq\\
&\qquad H^{5+i}(K(\pi_1(X),3+i), \pi_2(X))
\end{aligned}
\]

\uncertain{$S^1\wedge$} $K(\pi_1(X),1) \longrightarrow K(\pi_1(X),2)$

\begin{tikzcd}[column sep=small, row sep=small]
Y' \arrow[d] & \\
K(\pi_1(X),2) \arrow[r, "k_1'"] & K(\pi_2(X),4)
\end{tikzcd}

\begin{tikzcd}[column sep=small, row sep=small]
X = \Omega Y \arrow[d] & \\
K(\pi_1(X),1) \arrow[r, "k_1"] & K(\pi_2(X),3)
\end{tikzcd}

\page{83}
donc
\[
\Bigl(\frac{df}{f}\Bigr)^2 = \frac{(df)^2}{f^2} = \Bigl(\frac{f'}{f}\Bigr)^2 dt^2
\]
On a a) b) c) d) e) … On prend donc
\[
E(f,g) = \Bigl(\frac{dg}{g}\Bigr)^2 \Big/ \Bigl(\frac{df}{f}\Bigr)^2 = \frac{\bigl(\frac{dg}{df}\bigr)^2}{(g/f)^2}
\]

\struck{$\mathcal{E}_X/\mathbb{G}_a(\mathbb{C})$}

\page{85}
\struck{§} \uncertain{Schémas}.

\subsection*{1.1. Données équivalentes}

\begin{itemize}
\item[a)] fibré en droites projectives $P$ sur $X$ (fibré de Br.-Sev. de rang 1) \quad $f : P \to X$
\item[b)] « \struck{\ill{}} \uncertain{Module} loc. libre \add{de rang 2 $E$} mod. \uncertain{tensorisation} par un module inversible » sur $X$
\item[c)] Algèbre d'Azumaya \add{$M$} de rang 4
\item[d)] Algèbre de Lie $\mathfrak{g}$ forme \add{(tordue)} de $\mathfrak{gp}(1)_X$
\item[d')] \ldots\ $\mathfrak{g}'$ \ldots\ $\mathfrak{sl}(1)_X$
\item[e)] forme $G$ de $\mathrm{GP}(1)_X$
\item[e')] \ldots\ $G'$ de $\mathrm{SL}(2)_X$
\end{itemize}
\marginal{en regard de d) et d') : si 2 inv. sur $X$}
\note{les traits horizontaux de d') et e') sont ses « idem » ; l'indice de $\mathfrak{sl}(1)_X$ et de $\mathrm{SL}(2)_X$ est \uncertain{incertain} (« 1 » et « 2 » se confondent ici).}

Si on a $P$, on en déduit
\[
E = f_*(L),
\]
$L$ module inv. sur $P$ de degré 1 sur \add{les fibres} ($L$ défini loc. \add{sur $X$}, à \uncertain{tensorisation} près par $f^*(\mathcal{L})$, $\mathcal{L}$ inv. sur $X$),
\[
M = \check{E}\otimes E \simeq \underline{\mathrm{End}}(E) \simeq E\otimes E\otimes(\textstyle\bigwedge^2 E)^{-1}
\]
ext. de $\underline{O}_X\ (\simeq \bigwedge^2E\otimes(\bigwedge^2E)^{-1})$ par $\Gamma^2E\otimes(\bigwedge^2E)^{-1}$

ext. de $\underbrace{\mathrm{Sym}^2(E)\otimes(\bigwedge^2E)^{-1}}_{f_*(L^{\otimes 2})\otimes(\bigwedge^2E)^{-1} \simeq f_*(L^{\otimes 2}\otimes\bigwedge^2E^{-1})}$ par $\underline{O}_X\ (\simeq \bigwedge^2E\otimes(\bigwedge^2E)^{-1})$, où $L^{\otimes 2}\otimes\bigwedge^2E^{-1} = \underline{t}_{P/X}$.

\[
M \overset{?}{\simeq} f_*(\mathrm{P}^1_{P/X}(\underline{t}_{P/X}))\quad \text{\uncertain{addition}\,?}
\]

$\mathfrak{g} = M/\underline{O}_X$ avec crochet déduit du crochet \add{$(xy - yx)$} dans $M$ [\emph{NB} $\underline{O}_X$ est dans le centre]
\[
\simeq f_*(\underline{t}_{P/X})
\]
[avec crochet provenant des crochets dans $\underline{t}_{P/X}$]

$\mathfrak{g}' \subset M^{\mathrm{tr}=0}$\note{exposant « $\mathrm{tr}=0$ » \uncertain{douteux}} sous-algèbre des éléments de trace nulle ($\mathfrak{g}' \to \mathfrak{g}$ iso si 2 inv. sur $X$)

$\|$

$\mathfrak{g}' \subset \check{E}\otimes E$ sous-algèbre des \uncertain{éléments} de trace nulle

$G = \mathrm{Aut}_X P$, \quad $G' = \widetilde{G}$ (rev. univ. de Chevalley) $= \underline{\mathrm{SL}}(E)$ \struck{Aut \ill{}}

$P = \mathbb{P}(E) =$ schéma des Idéaux à gauche loc. \uncertain{facteurs directs} qui sont « \uncertain{spéciaux} » \struck{i.e. « \ill{} »} \ldots\ \struck{\ill{}}

$=$ schéma des sous-alg. de Borel de $\mathfrak{g}$ \add{$/\ \mathfrak{g}'$}

$=$ \ldots\ gr. de Borel de $G$ \add{$/\ G'$}

\marginal{$G \simeq \underline{\mathrm{Aut}}(E)/\mathbb{G}_{m,X} \simeq \underline{\mathrm{Aut}}(M) \simeq \underline{\mathrm{Aut}}(\mathfrak{g}) \simeq \underline{\mathrm{Aut}}(\mathfrak{g}')$ — ! attention à la car. (notamment car. 2 !)}

\marginal{explication [\ill{}] : $\mathrm{Sym}^2E\otimes\mathrm{Sym}^2E \to \mathrm{Sym}^2E\otimes\bigwedge^2E$, $\Gamma^2E\otimes\Gamma^2E \to \Gamma^2E\otimes\bigwedge^2E$ ; \struck{\ill{}} formes de Killing $\mathrm{Sym}^2E\otimes\mathrm{Sym}^2E \to (\bigwedge^2E)^{\otimes 2}$, $\Gamma^2E\otimes\Gamma^2E\to\ldots$}

\page{86}
$E =$ module d'une « rep. \uncertain{vraie} » de \struck{\ill{}} $M$, de $\mathfrak{g}$, de $\mathfrak{g}'$, de $G$, de $G'$.

$\mathfrak{g} = \underline{\mathrm{Lie}}\,G$

$\mathfrak{g}' = \mathrm{Lie}\,G'$

\subsection{1.2. Données équivalentes}

\begin{itemize}
\item[a)] fibré \struck{\ill{}} \add{en droites} proj. $P$ sur $X$, avec section unique $\sigma$
\item[b)] Structure d'extension
\[
0 \to \Omega \to \mathcal{L} \to \underline{O}_X \to 0,
\]
$\Omega$ inversible
\item[c)] Algèbre d'Azumaya de rang 4 $M$, avec idéal \uncertain{à gauche} « spécial » \struck{\ill{}} $\mathfrak{a}$
\item[d)] $\mathfrak{g}$ avec ss.-alg. de Borel $\mathcal{L}$
\item[d')] $\mathfrak{g}'$ \ldots\ $\mathcal{L}'$
\item[e)] $G$ \ldots\ gr. \ldots\ $B$
\item[e')] $G'$ \ldots\ gr. \ldots\ $B'$ \struck{\uncertain{tel qu'il}}
\item[f)] Une \struck{\ill{}} $\underline{O}_X$-Algèbre augmentée $\mathcal{B}_2$ \struck{\ill{}} \uncertain{nilpotente} d'\uncertain{échelon} 2, de « rang » 2 \add{à $\underline{O}_X$} et associée $\underline{L} = \underline{L}'$ faisceau inv. de degré relatif 1 rigidifié par $\sigma^*(\underline{L}) \simeq \underline{O}_X$. Alors
\end{itemize}
\note{l'alinéa f) est écrit plus serré et en partie par-dessus e') ; sa lecture est \uncertain{incertaine} et le passage de f) à « Alors $\mathcal{L} = f_*(\ldots)$ » ne se suit pas.}

$\mathcal{L} = f_*(\underline{L})$, avec structure d'ex. \uncertain{évidente}
\[
\simeq f_*(\underline{L}/\mathcal{J}\underline{L}) \qquad (\mathcal{J} = \text{Idéal de } \sigma(X) \text{ dans } P)
\]
\[
\simeq \sigma^*(\mathrm{P}^1(\underline{L}))
\]
$P = \mathbb{P}(\mathcal{L})$ avec section \uncertain{provenant} des \uncertain{quotients} inv. $\underline{O}_X$ de $\mathcal{L}$

$=$ \uncertain{classifiant} \add{\uncertain{proj.}} \ill{} des torseurs sous $W(\Omega)$ défini par l'extension

$M = \underline{\mathrm{End}}\,\mathcal{L}$, avec idéal annulateur de $\Omega$
\marginal{(données de \uncertain{splitting})}

\[
\mathfrak{g} = M/\underline{O}_X \simeq \mathrm{Sym}^2(\mathcal{L})\otimes\Omega^{-1} \qquad (\textit{NB}\ \det\mathcal{L}\simeq\Omega)
\]
Or $\mathrm{Sym}^2\mathcal{L}$ est filtré par \add{(décr.)} $\underline{O}_X, \Omega, \Omega^2$, donc $\mathrm{Sym}^2\mathcal{L}\otimes\Omega^{-1}$ filtré par \uncertain{(\ldots)}
\[
\begin{array}{ccc}
\Omega & \underline{O}_X & \Omega^{-1}\\
\| & \| & \|\\
\mathrm{Fil}_1 & \mathrm{Fil}_2/\mathrm{Fil}_1 & \mathrm{Fil}_3/\mathrm{Fil}_2\\
\| & \| & \|\\
\mathrm{gr}_1 & \mathrm{gr}_2 & \mathrm{gr}_3
\end{array}
\]
\marginal{expliquer aussi \uncertain{par} $\mathfrak{g}$, \uncertain{$\mathrm{P}^1(\mathcal{L})$} \ldots\ $0 \to \Omega \to \mathcal{L} \to \underline{O}_X \to 0$ \ldots\ ($\Omega = \mathcal{L}\ldots$)}

où $\mathrm{Fil}_2 \simeq \Omega^{-1}\otimes(\underbrace{\mathrm{Ker}\,\mathrm{Sym}^2\mathcal{L} \to \mathrm{Sym}^2 O}_{\simeq\ \Omega\otimes\mathcal{L}}) \simeq \mathcal{L}$

donc on a
\[
\underbrace{\Omega \subset \overbrace{\mathcal{L}}^{\underline{O}_X} \subset \mathfrak{g}}_{}
\qquad \text{(gradués : } \Omega,\ \underline{O}_X,\ \Omega^{-1})
\]
et $\mathcal{L}$ est une ss-alg. de Borel ($\Omega$ son prod. \uncertain{unipotent}) $\{$\emph{NB} pour forme de Killing, $\Omega$ et $\mathcal{L}$ orth. l'un de l'autre$\}$

\marginal{\emph{NB} $\mathfrak{g}/\Omega \simeq \mathcal{L}^{\vee}$ (vérifier) car Killing \ldots\ $\mathcal{L}\otimes\mathcal{L} \to \ldots$ \uncertain{non dégénérée}}

On n'explicite pas $(\mathfrak{g}', \mathcal{L}')$ ($\simeq(\mathfrak{g}, \mathcal{L})$ si 2 inv.), \uncertain{non plus} $(G,B)$, $(G',B')$ \ldots

$G' = \underline{\mathrm{SL}}_{\underline{O}_S}\mathcal{L}$, $B'$ \add{$\subset G'$} $=$ stabilisateur de $\Omega$

$G, B$ déduits par passage au quotient par centre \ldots

\marginal{plus bas}
[\uncertain{1.4}. \uncertain{Supposons} $X/S$, et considérons des \emph{connexions} \add{rel. $S$} sur les objets du type envisagé dans 1.1, mais auxquels on a ajouté \add{la donnée de} \uncertain{l'un} des structures 1.2. \uncertain{Ainsi}, on considère des couples \ldots
\[
(P, \sigma, c)
\]
\[
\begin{cases}
P \text{ fibré en droites projectives sur } X\\
\sigma \text{ section}\\
c \text{ connexion de } P \text{ rel. à } S
\end{cases}
\]
(mais où on \emph{suppose} \emph{pas} \ldots\ \uncertain{plat})

Il nous sera particulièrement commode de l'expliciter en termes de $(\mathcal{L},\Omega)$ i.e.
\[
0 \to \Omega \to \mathcal{L} \to \underline{O}_X \to 0
\]
(qui donne $(P,\sigma)$, \add{\uncertain{donc} $(\mathfrak{g},\mathcal{L})$, comme expliqué}, et d'une connexion de $\mathfrak{g}$. \struck{On notera que \ldots}]

\page{87}
\struck{\uncertain{Splittings}}

\subsection{1.3.}

fibré \uncertain{en droites} projectives sur $X$ avec \emph{deux} sections \emph{disjointes} $\sigma, \sigma'$

$\updownarrow$

fibré inversible $\Omega$ \quad ($P = \mathbb{P}(\Omega + \underline{O}_X) \simeq \overline{W(\Omega)}$ clôture projective \add{de $W(\Omega)$}, avec $\sigma =$ section à l'$\infty$, $\sigma' =$ section nulle)

$\updownarrow$

[$M$ somme directe de deux Idéaux \uncertain{spéciaux} \ldots]

$\updownarrow$

$\mathfrak{g}$ \add{\uncertain{munie}} \struck{\ill{}} \add{\uncertain{d'une sous-alg.} $\mathcal{L}$} \struck{\ill{} \uncertain{de Cartan}}, \uncertain{rang} 1 (cf. \uncertain{SGA 3}), soit d'une graduation en \ill{}
\[
\mathfrak{g} \simeq \underset{\substack{\|\\ \mathcal{L}_{-}}}{\Omega} + \underset{\substack{\|\\ \underline{O}_X}}{\underline{O}^0_X} + \Omega^{-1},
\qquad \mathcal{L} = \Omega + \underline{O}^0_X
\]
\[
\begin{aligned}
[\omega_1, \lambda] &= -[\lambda, \omega_1] = -\lambda\omega_1\\
[\omega_{-1}, \lambda] &= -[\lambda, \omega_{-1}] = \lambda\omega_{-1}\\
[\omega_1, \omega_{-1}] &= -[\omega_{-1}, \omega_1] = 2\omega_1\omega_{-1}
\end{aligned}
\]
soit $\mathfrak{g}$ avec $\mathfrak{t}\subset\mathfrak{g}$ \uncertain{sous-alg.} de Cartan, et le choix d'une des deux sous-algèbres \uncertain{nilpotentes} \add{$\mathfrak{t}_{\mathfrak{g}}$} normalisées par $\mathfrak{t}$, $\Omega, \Omega'$ \ldots\ \uncertain{ce qui fixe} l'iso $\mathfrak{t} \simeq \underline{O}_X$ ($\simeq \underline{\mathrm{Hom}}_{\underline{O}_S}(\Omega,\Omega)$)
\marginal{\struck{\ill{}} $\mathfrak{g}'$ \ldots\ sauf \uncertain{car.} $\neq 2, 3$ \ldots\ \uncertain{modulo} \ill{}}

$\updownarrow$ \quad $\mathfrak{g}'$ avec données analogues (Mais dans la table de multiplication, dans ce cas, il faut \ldots
\[
-2\lambda\omega_1,\quad 2\lambda\omega_{-1},\quad \omega_1\omega_{-1}\ )
\]

$\updownarrow$ \quad $G$ avec couple de Killing

$\updownarrow$ \quad $G'$ avec couple de Killing
\note{« couple de Killing » est une lecture \uncertain{probable} des deux dernières lignes.}

\page{88}
\subsection{1.4 suite}

Choisissons d'abord un splitting de $\mathfrak{g}$, \add{\uncertain{rel.} $\mathcal{L}$ i.e. un}
\[
\mathfrak{g} \simeq \Omega + \underset{\substack{\|\\ \underline{O}_X}}{\Omega^0} + \Omega^{-1}
\]
d'où, posant $\omega = \underline{\Omega}^1_{X/S}$,
\[
\mathfrak{g}\otimes\omega \simeq \Omega\otimes\omega + \omega + \Omega^{-1}\otimes\omega
\]
\struck{\ill{}}

Choisissons-en loc. une section $\delta$ de $\Omega$ comme base. Une connexion (sans courbure) $c_0$ en \uncertain{obtenue} \uncertain{par}
\[
\mathfrak{g} \xrightarrow{\ c_0\ } \mathfrak{g}\otimes\omega
\]
\[
c_0(\delta) = c_0(1) = c_0(\delta^{-1}) = 0
\]
i.e.
\[
c_0(\underbrace{\lambda\delta + \mu\cdot 1 + \nu\delta^{-1}}_{\substack{\updownarrow\\ (\lambda,\mu,\nu)}}) = \underbrace{\delta\,d\lambda + 1\,d\mu + \delta^{-1}d\nu}_{\substack{\updownarrow\\ (d\lambda,\,d\mu,\,d\nu)}}
\]
Une connexion générale s'obtient alors en \uncertain{ajoutant} une section de
\[
\mathrm{Der}(\mathfrak{g},\mathfrak{g})\otimes\omega \simeq \mathfrak{g}\otimes\omega
\]
(\emph{NB} toute dérivation de $\mathfrak{g}$ est intérieure !)

Soit $\varpi = \varpi_2 + \varpi_1 + \varpi_0 \in \Gamma\,\mathfrak{g}\otimes\omega$ i.e.
\[
\begin{aligned}
\varpi_2 &\in \Gamma(\Omega\otimes\omega)\\
\varpi_1 &\in \Gamma(\omega)\\
\varpi_0 &\in \Gamma(\Omega^{-1}\otimes\omega)
\end{aligned}
\]
on aura donc pour connexion $c_{\varpi}$ alors \uncertain{par} $c_0$ et $\varpi$
\[
c_{\varpi}(\underbrace{\lambda\delta + \mu\cdot 1 + \nu\delta^{-1}}_{(\lambda,\mu,\nu)}) \overset{?}{=} \text{\struck{\ill{}}}
\]

\page{89}
\[
\begin{aligned}
c_{\varpi}(\underset{\substack{\|\\ \omega_0}}{\lambda\delta}) &= \underbrace{\delta\,d\lambda}_{\text{poids } 2} + \lambda\bigl[\underset{\substack{\|\\ 0}}{[\varpi_2,\lambda\delta]} + [\varpi_1,\lambda\delta] + [\varpi_0,\lambda\delta]\bigr]\\
&\qquad + \underbrace{\varpi_1(\lambda\delta)}_{\text{poids } 2}\ \ 2\varpi_0(\lambda\delta)
\end{aligned}
\]
\note{la première ligne est surchargée ; l'indice de $\omega_0$ sous $\lambda\delta$ et le signe devant $2\varpi_0(\lambda\delta)$ sont \uncertain{incertains}.}

\[
\begin{aligned}
c_{\varpi}(\underset{\substack{\|\\ \lambda\delta}}{\omega_1}) &= \overset{2}{\delta\,d\lambda} + \omega_1(\overset{2}{\varpi_1} - \overset{1}{2\varpi_0})\\
c_{\varpi}(\underset{\substack{\|\\ \mu\cdot 1}}{\omega_0}) &= \overset{1}{d\mu} + \omega_0(-\overset{2}{\varpi_2} + \overset{0}{\varpi_0})\\
c_{\varpi}(\underset{\substack{\|\\ \nu\delta^{-1}}}{\omega_{-1}}) &= \overset{0}{\delta^{-1}d\nu} + \omega_{-1}(\overset{1}{2\varpi_2} - \overset{0}{\varpi_1})
\end{aligned}
\]
\note{les petits chiffres au-dessus des termes sont ses indices de filtration (« poids ») ; leur position exacte est \uncertain{approximative}.}

Encadré :
\[
\mathfrak{g} : \Omega,\ \underset{\substack{\|\\ \underline{O}_X}}{\Omega^0},\ \Omega^{-1}
\qquad\qquad
\mathfrak{g}\otimes\underline{\omega} : \Omega\otimes\underline{\omega},\ \underline{\omega},\ \Omega^{-1}\otimes\underline{\omega}
\]

On trouve la condition suivante (indépendante du choix des splittings \ldots !)

$\{$ \emph{La connexion $c$ satisfait} $c(\mathrm{Fil}_i\,\mathfrak{g}) \subset \mathrm{Fil}_{i+1}(\mathfrak{g}\otimes\underline{\omega})$

$\{$ par passage au quotient, il induit
\[
c_{i,i+1} : \mathrm{gr}_i\,\mathfrak{g} \to \mathrm{gr}_{i+1}(\mathfrak{g}\otimes\omega) \simeq \mathrm{gr}_{i+1}(\mathfrak{g})\otimes\omega
\]
qui est \emph{linéaire}, de façon précise
\[
\begin{aligned}
c_{1,2} &: \mathrm{gr}_1 \to \mathrm{gr}_2,\quad \Omega \to \omega && \text{i.e. } c_{12} \in \Gamma\,\Omega^{-1}\otimes\underline{\omega}\\
c_{2,3} &: \mathrm{gr}_2 \to \mathrm{gr}_3,\quad \underline{O}_X \to \Omega^{-1}\otimes\omega && \text{i.e. } c_{23} \in \Gamma\,\Omega^{-1}\otimes\underline{\omega}
\end{aligned}
\]
on a $[\,c_{12} = -2\varpi_0,\ c_{2,3} = \varpi_0$, donc$\,]$
\[
\boxed{c_{1,2} = -2\,c_{2,3}}
\]

\emph{NB} On \uncertain{déduit} que la connexion respecte la filtration ssi $\underset{\sim}{c_{2,3}} = 0$ (d'où $c_{1,2} = 0$) [cela signifie $\varpi_0 = 0$] \struck{\ill{}} \add{\uncertain{ce sont aussi}} celles qui respectent la Borel $\mathcal{L}\subset\mathfrak{g}$, ou encore \struck{\ill{}} qui respectent la section $\sigma$ de $P$, i.e. pour lesquelles $\sigma$ est horizontale. Elles sont déterminées \uncertain{modulo} section de $\mathcal{L}\otimes\underline{\omega}$.

\struck{\ill{}}

Re \emph{NB} On pourrait aussi considérer
\[
c^{\Omega} : \underset{\substack{\wr\\ \mathcal{L}'\simeq\mathcal{L}\otimes\Omega^{-1}}}{\mathfrak{g}/\Omega} \longrightarrow \Omega^{-1}\otimes\underline{\omega}
\]
induit par $c$, \ill{} \uncertain{(bilinéaire)} sur $\Omega^0 \simeq \underline{O}_X \subset \mathcal{L}\otimes\Omega^{-1}$ ; sa connaissance est \emph{équivalente} à celle de $c_{\Omega}$ si 2 inv., etc. dans le cas général (\uncertain{précisément} car. 2 ? C'est \uncertain{peut-être} plus précis \ldots)

\emph{NB} Les connexions sur $\mathfrak{g}$ compatibles avec crochet \struck{\ill{}} (si 2 inv. --- sinon il faut parler des connexions sur $P$, ou $G$ etc.) forment un torseur sous $\mathfrak{g}\otimes\underline{\omega}$. Or
\[
\text{on a}\quad \mathfrak{g}\otimes\underline{\omega} \to \Omega^{-1}\otimes\underline{\omega}
\]
d'où un torseur sous $\Omega^{-1}\otimes\underline{\omega}$. \emph{Le dernier est splitté canoniquement}, i.e. on a
\[
\underline{\mathrm{Conn}} \xrightarrow{\ \varpi_0\ } \Omega^{-1}\otimes\omega
\]
compatible avec $\varpi_0$, et si $c$ section de $\underline{\mathrm{Conn}}$, $\varpi_0(c)$ détermine les homs $\mathrm{gr}_i(\mathfrak{g}) \to \mathrm{gr}_{i+1}(\mathfrak{g}\otimes\underline{\omega})$
\[
\begin{aligned}
\Omega = \mathrm{gr}_1(\mathfrak{g}) &\xrightarrow[-2\varpi_0]{} \omega \simeq \mathrm{gr}_2(\mathfrak{g}\otimes\underline{\omega})\\
\Omega^0 \simeq \underline{O}_X = \mathrm{gr}_2\,\mathfrak{g} &\xrightarrow[\varpi_0]{} \Omega^{-1}\otimes\underline{\omega} = \mathrm{gr}_3(\mathfrak{g}\otimes\underline{\omega})
\end{aligned}
\]
Si on considère les $c$ pour lesquelles $\varpi_0(c) \in \Gamma\,\underset{\substack{\|\\ \mathrm{Hom}(\Omega,\underline{\omega})}}{\Omega^{-1}\otimes\underline{\omega}}$ est donnée $i$, on trouve une restriction du gr. struct. de $\underline{\mathrm{Conn}}$ : $\mathcal{L}\otimes\underline{\omega}$, d'où torseur $\underline{\mathrm{Conn}}_i$ sous $\mathcal{L}\otimes\underline{\omega}$. Or on a
\[
\mathcal{L}\otimes\underline{\omega} \to \underline{\omega}
\]
d'où un torseur \uncertain{induit} sous $\underline{\omega}$ : les conn. sur $\mathfrak{g}$ compatibles avec $i$, \emph{modulo} \add{op. d'} \uncertain{ajouter} une section de $\Omega\otimes\underline{\omega} \subset \mathcal{L}\otimes\underline{\omega} \subset \mathfrak{g}\otimes\underline{\omega}$. C'est aussi \struck{(\ill{})} le torseur des restrictions à $\Omega$ des conn. compatibles avec $i$, ou encore des homs \uncertain{additifs} op. diff. d'ordre 1 (rel. $S$)
\[
\Omega \xrightarrow{\ c_{\Omega}\ } \mathcal{L}\otimes\omega
\]
telles que la composée \struck{\ill{}} $\Omega \to \mathcal{L}\otimes\omega \to \omega$ soit linéaire (*) (et égale à $-2i$), et qu'on ait
\[
(**)\quad c_{\Omega}(f\omega) - f\,c_{\Omega}(\omega) = \omega\otimes df \qquad \omega\in\Gamma\,\Omega,\ f\in\Gamma\,\underline{O}_X
\]
\marginal{TSVP}

\page{90}
La donnée d'un tel $c_{\Omega}$ équivaut à une restriction du groupe structural de $\underline{\mathrm{Conn}}_i$ à $\Omega\otimes\underline{\omega}$, ou encore (indépendant \struck{\uncertain{de la donnée de}} \add{de} $i$) la donnée d'un $c_{\Omega}$ pour lequel (*) est linéaire et (**) est vraie, équivaut à une restriction du groupe structural de $\mathfrak{g}\otimes\omega$ à $\underline{\mathrm{Conn}}$ \ldots\ $\Omega\otimes\underline{\omega}$ \ldots

On s'intéresse au cas où $\underline{\mathrm{Conn}}_i$ est \uncertain{trivial}, i.e. $c_{\Omega}$ est tel que le composé
\[
(-2i) : \Omega \longrightarrow \mathcal{L}\otimes\underline{\omega} \to \underline{\omega}
\]
\struck{(ce qui exige que 2 inv.)} \struck{\ill{}}

soit un \emph{iso}. Alors les triples
\[
(\mathfrak{g}, \mathcal{L}, c_{\Omega}) \quad (\text{où, ce qui revient au même,}\ (\mathcal{L}\supset\Omega,\ c_{\Omega}))
\]
sont définis, à iso unique près, par $X/S$ (et \uncertain{avant} \uncertain{tout} $\underline{\omega} = \underline{\Omega}^1_{X/S}$ est inversible) comme
\[
\mathcal{L} \simeq \mathrm{P}^1_{X/S}(\underline{\omega})\otimes\underline{\omega}^{-1}
\qquad (\text{d'où } \mathcal{L}\otimes\underline{\omega} \simeq \mathrm{P}^1_{X/S}(\underline{\omega}))
\]
et
\[
c_{\Omega} : \Omega \simeq \underline{\omega} \longrightarrow \mathcal{L}\otimes\underline{\omega} \simeq \mathrm{P}^1_{X/S}(\underline{\omega})
\]
n'est autre que $d^1_{X/S,\underline{\omega}}$ \ldots

\page{91}
Interprétation de $c_{23}$ comme
\[
c_{2,3} : \Omega \longrightarrow \underline{\omega} = \Omega^1_{X/S}
\]
\ldots\ une fois qu'on a \uncertain{noté}, compte tenu de
\[
\Omega \simeq \sigma^*(\underline{\Omega}^1_{P/X})
\]
la division relative à la connexion $c$ de la section $\sigma$ ; on dit que $\sigma$ est loc. (\ldots\ i.e.\ $\mathcal{L}\subset\mathfrak{g}$) est \emph{nulle part stationnaire} pour $c$, si $c_{2,3}$ est « partout non nulle ». \struck{\uncertain{Appliquer} \ldots} \add{\uncertain{Il faut} \ldots} dire \uncertain{local} du $\omega$.

On suppose par la suite qu'il en est ainsi, \emph{i.e.} que $\omega$ est \uncertain{inversible}, \uncertain{et} \uncertain{supposons} par $c_{2,3} : \Omega \to \underline{\omega}$ est un iso. On \uncertain{appellera} les connexions « \emph{transverses} à $\sigma$ (\uncertain{ou} à $\mathcal{L}\subset\mathfrak{g}$) ». On identifie $\Omega\simeq\underline{\omega}$ par $c_{23}$, et de façon précise, \uncertain{on} suppose donné un iso
\[
\Omega \xrightarrow[\ \sim\ ]{\ i\ } \underline{\omega} = \Omega^1_{X/S}
\]
et on se borne aux conn. pour lesquelles $c_{2,3} = i$. \struck{\ill{}} \add{\uncertain{Les} \ill{}} sections d'un torseur sous $\mathcal{L}\otimes\underline{\omega}$ [soit $\underline{\mathrm{Conn}}(\mathfrak{g}, i)$ \add{\uncertain{transverses}}].

Considérons la restriction d'une connexion $c$ à $\Omega \simeq \underline{\omega} = \mathrm{Fil}_1\mathfrak{g}$
\[
c|\underset{\substack{\|\\ \mathrm{Fil}_1\mathfrak{g}}}{\Omega} : \Omega \longrightarrow \underbrace{\mathcal{L}\otimes\omega}_{\Omega\otimes\omega,\ \omega} \ldots
\]
\note{la page s'arrête sur cette formule inachevée.}

\page{92}
Je sais que son composé avec $\mathcal{L}\otimes\underline{\omega} \to \underline{\omega}$ est $-2i\ldots$. \struck{Considérons} Comme \struck{$c|\Omega$} \add{$c|\underline{\omega}$} est op. diff. d'ordre $\leq 1$, on trouve qu'il se factorise en
\[
\underline{\Omega} \xrightarrow[\substack{\text{op. diff.}\\ \text{con.}\\ \text{d'ordre } 1}]{\ d^1_{\underline{\omega}}\ } \mathrm{P}^1_{X/S}(\underline{\Omega}) \xrightarrow[\text{lin.}]{\ \gamma\ } \mathcal{L}\otimes\underline{\omega}
\]
et comme le composé \struck{$c|\underline{\omega}$}
\[
\underline{\Omega} \xrightarrow{\ c|\underline{\omega}\ } \mathcal{L}\otimes\underline{\omega} \to \underline{\omega}
\]
est linéaire (de valeur $-2i$) on trouve que $\gamma$ définit une \ill{} linéaire d'extensions
\[
\begin{array}{ccccccccc}
0 \to & \Omega\otimes\omega & \to & \mathrm{P}^1_{X/S}(\Omega) & \to & \Omega & \to & 0\\
& \downarrow{\scriptstyle \gamma_0 = ?} & & \downarrow{\scriptstyle \gamma} & & \downarrow{\scriptstyle -2i} & &\\
0 \to & \Omega\otimes\omega & \to & \mathcal{L}\otimes\omega & \to & \omega & \to & 0
\end{array}
\]
On va déterminer $\gamma_0$, quels que soient les splittings $i$ et section $\delta_0$ de $\Omega$, en termes de $\varpi_0 = i$, $\varpi_1$, $\varpi_2$ ci-dessus, \struck{\ill{}} \add{\uncertain{En fait} \ill{}} \struck{\uncertain{$\mathrm{P}^1_{X/S}(\Omega) \simeq \Gamma\,\mathrm{P}^1_{X/S}(\underline{O}_X)$} \ill{}} \ldots\ \uncertain{op. diff.} \uncertain{loc. libre} \ldots\ $\Omega$, \ldots

\note{les deux dernières lignes de l'alinéa sont serrées, en partie biffées, et lues pour l'essentiel \ill{}.}

avec
\[
\gamma(\underset{\substack{\cap\\ \Gamma\Omega}}{\omega_1}\otimes \underset{\substack{\cap\\ \Gamma\underline{O}_X}}{df}) = c(f\omega_1) - f\,c(\omega_1)
\]
Si $\omega_1 = \lambda\delta$, on trouve donc
\[
\begin{aligned}
\gamma(\lambda\delta\otimes df) &= \delta\,[d(f\lambda) - f\,d(\lambda)]\\
&= \delta\,[\lambda\,df] = \lambda\delta\otimes df
\end{aligned}
\]
donc $\gamma = \mathrm{id}$. Donc

\emph{Remarque.} Si $c$ connexion \emph{transverse} \add{et 2 inv. sur $X$}, alors \uncertain{$\gamma_0$} est un \emph{iso} d'extensions de $\mathrm{P}^1_{X/S}(\Omega)$ \ldots
\note{la phrase se poursuit à la page suivante.}

\page{93}
et $\mathcal{L}\otimes\underline{\omega}$, induisant $-2i$ et $\mathrm{id}$. Donc un iso
\[
\mathcal{L} \xrightarrow[\gamma^{-1}\otimes\mathrm{id}_{\omega^{-1}}]{\ \sim\ } \mathrm{P}^1_{X/S}(\Omega)\otimes\omega^{-1} \simeq \mathrm{P}^1_{X/S}(\omega)\otimes\omega^{-1}
\]
\struck{Corr.}

induisant un diagramme d'extensions

\begin{tikzcd}[column sep=small, row sep=small]
0 \arrow[r] & \underline{\omega} \arrow[r] \arrow[d, "\mathrm{id}"] & \mathcal{L} \arrow[r] \arrow[d] & \underline{O}_X \arrow[r] \arrow[d, "-2\,\mathrm{id}"] & 0\\
0 \arrow[r] & \underline{\omega} \arrow[r] & \mathrm{P}^1_{X/S}(\omega)\otimes\omega^{-1} \arrow[r] & \underline{O}_X \arrow[r] & 0
\end{tikzcd}
\marginal{(R) en regard du diagramme ; $-2\,\mathrm{id}$ est cerclé.}

ou encore, en multipliant par $-\frac{1}{2}$

\begin{tikzcd}[column sep=small, row sep=small]
0 \arrow[r] & \underline{\omega} \arrow[r] \arrow[d, "-\frac{1}{2}\mathrm{id}_{\underline{\omega}}"] & \mathcal{L} \arrow[r] \arrow[d] & \underline{O}_X \arrow[r] \arrow[d, "\mathrm{id}_{\underline{O}_X}"] & 0\\
0 \arrow[r] & \underline{\omega} \arrow[r] & \mathrm{P}^1_{X/S}(\omega)\otimes\omega^{-1} \arrow[r] & \underline{O}_X \arrow[r] & 0
\end{tikzcd}
\marginal{(RR) en regard du second diagramme ; $-\frac{1}{2}\mathrm{id}_{\underline{\omega}}$ est cerclé.}

(Dans le torseur \uncertain{sous} $\underline{\omega} \simeq \Omega$ défini par l'ext. $\mathcal{L}$ est déduit du torseur défini \add{\uncertain{par} $\underline{\omega}$} l'ext. $\mathrm{P}^1_{X/S}(\underline{\omega})\otimes\underline{\omega}^{-1}$ de $\underline{O}_X$ par $\underline{\omega}$ (\uncertain{ou} $\mathrm{P}^1_{X/S}(\omega)$ de $\underline{\omega}$ par $\underline{\omega}^2$) par \emph{multiplication par $-2$}.)

Utilisant cette identification, $c|\Omega \simeq \underline{\omega}$ s'identifie à l'op. diff. universel
\[
c\,|\,\underset{\substack{\|\\ \mathrm{Fil}_1\mathfrak{g}\\ \|\\ \mathcal{L}^{\perp}}}{\Omega \simeq \underline{\omega}}
\quad \underset{d^1_{\underline{\omega},X/S}}{=} \quad
\underset{\substack{\wr\\ \omega}}{\Omega} \xrightarrow{\ c\ } \mathcal{L}\otimes\omega \simeq \mathrm{P}^1(\omega)
\]
\note{l'exposant de $\mathcal{L}$ sous $\mathrm{Fil}_1\mathfrak{g}$ est lu « orth. » \uncertain{(orthogonal)}.}

\emph{NB} On utilise \struck{d'ici} \add{ici} ces mêmes identifications et on se borne aux conn. de $\mathfrak{g}$, \struck{\ill{}} dont la restriction à $\Omega \simeq \underline{\omega} = \mathcal{L}$ est la $d^1_{\underline{\omega},X/S}$. Elles sont déterminées \add{\uncertain{donc}} \uncertain{modulo} \struck{\uncertain{faisant}} \add{sections} \struck{\uncertain{dans}} de $\Omega\otimes\omega \simeq \underline{\omega}^{\otimes 2}$.

\page{94}
\emph{Donc} pour la suite on se donne d'abord l'extension
\[
0 \to \Omega \to \mathcal{L} \to \underline{O}_X \to 0
\]
($\Omega$ inversible), puis l'op. diff. d'ordre $\leq 1$
\[
\Omega \xrightarrow{\ c_{\Omega}\ } \mathcal{L}\otimes\underline{\omega} \qquad \underline{\omega} = \Omega^1_{X/S}
\]
tel que
\begin{itemize}
\item[a)] le composé avec $\mathcal{L}\otimes\underline{\omega} \to \underline{\omega}$ soit \uncertain{un}
$\Omega \to \underline{\omega}$ linéaire et un \emph{iso} (qui \uncertain{servira} à identifier $\Omega \simeq \underline{\omega}$ \uncertain{puis} \uncertain{utiliser} \uncertain{plus} \ldots\ $-2i$, $i : \Omega \simeq \ldots$ pour identifier).
\item[b)] \struck{\ill{}} On a (si $\omega\in\Gamma\Omega$, $f\in\Gamma\underline{O}_X$)
\[
c_{\Omega}(f\omega) - f\,c_{\Omega}(\omega) = \omega\otimes df
\]
($\omega$ fixe, c'est \uncertain{bien} \uncertain{linéaire} $\Omega \to \Omega\otimes\omega$)
\end{itemize}
\marginal{cela \uncertain{équivaut} \uncertain{vraiment} \uncertain{ici} \ldots}

[Les conditions définissent le splitting de l'ext. $\mathcal{L}$ et l'op. diff. $c_{\Omega}$ \uncertain{à iso} unique près --- et ce pour types \ill{} \ldots\ si $\underline{\omega} = \Omega^1_{X/S}$ est bien inversible)

On forme $\mathfrak{g}$, et on considère les connexions [\add{sur $\mathfrak{g}$}] d'alg. de Lie $\underline{\mathrm{Conn}}$ sur $\mathfrak{g}$ qui prolongent $c_{\Omega}$. Elles forment les sections d'un torseur sous $\Omega^2\otimes\underline{\omega}$ $\simeq \underline{\omega}^{\otimes 2}$. \struck{On \ill{} \ldots\ qu'elles \ldots}

\marginal{Le donnée de $c_{\underline{\omega}}$ revient à \ill{} d'un iso. des \ill{} de $\Omega$ \ldots\ $\sigma(X)$ de $X\times_S X$ \ldots\ \uncertain{voisinage} \ldots\ $\sigma(X) \ldots$ \ldots\ \uncertain{voisinage} \ldots\ $\ldots$ \ill{} \ldots}
\note{note marginale écrite verticalement, dans un cadre, en partie biffée ; elle n'est lue que par fragments.}

Faire le lien avec la description $(P,\sigma,c)$ et les deux familles
\[
\underline{O}_{\hat{P}} \simeq \mathrm{P}^{\infty}_{X/S}
\]
(si $X/S$ diff\supplied{érentiellemen}t lisse \ldots).

\page{95}
\subsection{1.5. Sections horizontales de $P$}

Pour quelles connexions $\mathfrak{g}$ \uncertain{admissibles} sur $X/S$ (avec $\underline{\omega}_{X/S}$ \uncertain{inversible}) y a-t-il des sections horizontales de $P$, \struck{et} quelles sont-elles ? On a étudié \uncertain{déjà} les sections horizontales de $P$ là où elles sont disjointes de la section donnée $\sigma$. Elles correspondent donc à une décomposition de $\mathfrak{g}$ en somme directe, \ldots\ de $\mathfrak{g}$ en
\[
\mathfrak{g} \simeq \Omega \oplus \Omega^0 \oplus \Omega^{-1} \qquad (\Omega\simeq\omega)
\]
\struck{Si le \ill{}} Une connexion $c$ \uncertain{admet} la section $\sigma'$ horizontale \struck{$c(\Omega^{-1})\subset$} ssi $c(\Omega^{-1}) \subset (\Omega^0 + \Omega^{-1})\otimes\underline{\omega}$ \add{ou $c(\Omega^0)\subset(\Omega^0+\Omega^{-1})\otimes\underline{\omega}$}, \ldots\ comme \uncertain{on} \uncertain{le} \uncertain{trouve} \uncertain{en} \uncertain{termes} d'une \struck{section} \add{base} $\delta$ de $\Omega$, permettant de décrire $c$ par trois formes $\varpi_2, \varpi_1, \varpi_0$, sections de $\Omega\otimes\underline{\omega}$, $\Omega^0\otimes\underline{\omega} = \underline{\omega}$ et $\Omega^{-1}\otimes\underline{\omega}$ :
\[
-2\varpi_2 = 0
\]
\struck{\uncertain{donc si 2 inversible cela signifie} $\varpi_2 = 0$} et pour \ill{} obtenir \uncertain{ainsi} une \ill{} de connexions \uncertain{modulo} $\Gamma(\Omega\otimes\omega)$ et donc \ldots\ \struck{\uncertain{p. ex.}} \ill{} par sa restriction à $\Omega$ --- \uncertain{puis} \uncertain{on} obtient \uncertain{des} \ldots\ unique en \uncertain{convergeant} $c$ par \uncertain{adjonction} d'une $\omega_2\in\Gamma\,\Omega\otimes\underline{\omega}$.
\marginal{$2$ \uncertain{inv.} !}

Donc

\emph{Prop.} Pour \uncertain{le} splitting de \add{l'ext.} $\mathcal{L}$ ($\simeq \mathrm{P}^1_{X/S}(\underline{\omega})\otimes\underline{\omega}^{-1}$ dans le cas « canonique ») il y a une et une seule conn. \emph{admissible} sur $\mathfrak{g}$ (\uncertain{ou} $P\ldots$) pour laquelle la section correspondante de $P$ soit horizontale. Donc \struck{\uncertain{correspondance} \ldots\ \uncertain{bien}} \add{\ill{}} \ldots\ l'ensemble des sections de $P - \sigma(X)$ \struck{\ill{}} \add{\uncertain{affines}} \uncertain{à} \ldots\ \add{i.e. des splittings de $\mathrm{P}^1_{X/S}(\omega)$} sur $X$, \struck{\ill{}} \add{\uncertain{et}} des connexions \struck{\ill{}} \add{\uncertain{affines}} sur $X$, ayant comme \ldots\ \uncertain{que} celles qui admettent \uncertain{une} section hor. disjointe de $\sigma$ ;
\note{la fin de la page est très surchargée de biffures et d'ajouts interlinéaires ; la lecture n'en est que partielle.}

\page{96}
Quelles connexions \uncertain{sont} \uncertain{obtenues} \uncertain{ainsi} ? \struck{Choisissons (loc.) une uniformisante} \struck{$t$ de $X$, d'où une base $dt$ de $\omega\simeq\Omega$, et un splitting de $\mathrm{P}^1_{X/S}(\Omega) \simeq \mathrm{P}^1_{X/S}(\omega) \simeq \mathrm{P}^1(\underline{O}_X)$ (par le dual de la base $dt$), d'où une connexion \add{\uncertain{régulière}}, correspondante.} On va supposer $\mathfrak{g}$ splitté localement, avec base $\delta$ de $\Omega$, et une connexion d'origine $c_0$ donnée par
\[
\begin{aligned}
c_0(\underset{\substack{\|\\ \lambda\delta}}{\omega_1}) &= \delta\,d\lambda + \omega_1(-2\varpi_0) \qquad\qquad \varpi_0\in\Gamma\,\Omega^{-1}\otimes\omega\\
c_0(\underset{\substack{\|\\ \mu\cdot 1}}{\omega_0}) &= 1\,d\mu + \omega_0(\varpi_0)\\
c_0(\underset{\substack{\|\\ \nu\delta^{-1}}}{\omega_{-1}}) &= \delta^{-1}d\nu \quad \text{\struck{\ill{}}}
\end{aligned}
\]
Considérons une section $s$ de $\mathcal{L}$ s'envoyant sur la section $1_X$ de $\underline{O}_X$
\[
s = s_0 + \omega_1 \qquad (s_0 = (0,1))
\]
Soit
\[
\begin{cases}
\text{\struck{$\ldots$}} : \underline{O}_X \to \mathcal{L}\\
v_s : \Omega^{-1} \to \mathfrak{g}
\end{cases}
\qquad \text{i.e. } v_s \in \Gamma(\mathfrak{g}\otimes\Omega)
\]
\struck{\uncertain{les}} \add{\uncertain{les}} \uncertain{\ill{}} \uncertain{correspondants}. \struck{de} On a
\[
\begin{aligned}
s &= \exp\Theta_{\omega_1} s_0 \quad \text{\struck{$= s_0 + [-\omega_1, s_0] + [-\omega_1 \ldots$}}\\
v_s &= \exp\Theta_{\omega_1} v_{s_0} = v_{s_0} + \underbrace{[-\omega_1, v_{s_0}]}_{-2\omega_1 v_{s_0}} + \tfrac{1}{2}\underbrace{[-\omega_1,[-\omega_1, v_{s_0}]]}_{\omega_1^2 v_{s_0}}
\end{aligned}
\]
\marginal{pour $v_{s_0} = \ldots$}

\struck{d'où} la question \struck{\uncertain{est de}} \struck{$c(v_s)$ \uncertain{est}} de calculer la \uncertain{composante} \add{$\varpi_2(s) = \varphi(\omega_1)$} de $c(s)$ \uncertain{suivant} le \struck{\ill{}} \add{$\Omega\otimes\omega$} \struck{\uncertain{de la décomposition}}
\[
\mathfrak{g} \simeq \Omega + \text{\struck{$\exp\Theta_{\ldots}$}}\,\underline{O}_X\cdot s + v_s(\Omega^{-1})
\]
i.e. $(\exp\Theta_{\omega_1})\otimes\mathrm{id}_{\omega}(c(s))$ \struck{et la forme \ldots\ (pour \ldots)} \struck{$\varpi_2(s) = \varphi(\omega_1)$, \ldots}
\[
c(v_s) = \underbrace{\varpi_2(s)}_{\Omega\otimes\omega} + \text{\struck{\ill{}}}\ s\otimes\varpi_1(\ldots) + \underbrace{(v_s\ldots)(\varpi_0)}_{\Omega^{-1}\otimes\omega}
\]
\[
= \text{composante de } \exp\Theta_{+\omega_1}(c(s))
\]
\note{ce bloc est entouré, traversé de biffures et de flèches de renvoi ; sa lecture est en grande partie \uncertain{incertaine}.}

Sauf erreur, on trouve
\[
s = s_0 + \omega_1 \longmapsto c_0 + \Bigl(\delta\,d\bigl(\tfrac{\omega_1}{\delta}\delta^{-1}\bigr) + \omega_1^2\Bigr)
\]
Prenons $\delta$ de façon que son image dans $\omega$ soit $dt$, $t$ une uniformisante, on trouve sauf erreur de calcul
\[
s_0 + \lambda\,dt \longmapsto c_0 + \Bigl(\frac{d\lambda}{dt} + \lambda^2\Bigr)(dt)^2
\]
\marginal{facteur 2 ?}

Donc quand une conn. proj. sur $X/S$ est définie par \uncertain{l'opération} de la \struck{\ill{}} conn. proj. de réf. $c_0$) par une 2-forme $\varphi(dt)^2$, la question s'il y a \uncertain{possibilité} de trouver une section horizontale de $P$ disjointe de la section \struck{donnée} \add{$\sigma$} revient à celle de trouver des solutions $\lambda$ de l'équation
\[
\frac{d\lambda}{dt} + \lambda^2 = \varphi
\]
(équation diff. \uncertain{non} lin. du 1er ordre). On sait, si $X$ courbe lisse sur corps \uncertain{\ill{}} de car. nulle, qu'il y a au plus une solution ayant une valeur et une dérivée données en un pt ; et il y en a une \uncertain{formelle}, mais \uncertain{transcendante} (!), en \uncertain{\ill{}} générale --- les pôles de \add{\uncertain{l'unique}} $\lambda$ correspondant \struck{\ill{}} \add{\uncertain{seuls}} à des pts de rencontre de $\sigma$ (défini par $\lambda$) et de $\sigma'\ldots$

\subsection{1.6.}

Intrinsèquement, le $P$ canonique est $\simeq \mathbb{P}(\mathrm{P}^1_{X/S}(\underline{\omega}))$, où $\underline{\omega} = \Omega^1_{X/S}$. Donc une section de $P$ \struck{\ill{}} s'identifie à un couple
\[
(L, \sigma)
\]
$L$ Module inversible sur $X$, et
\[
\sigma : \mathrm{P}^1_{X/S}(\omega) \to L \quad \text{épi. lin.}
\]
Or un hom. \add{lin.} $\sigma : \mathrm{P}^1_{X/S}(\omega) \to L$ équivaut à \ldots
\note{la phrase se poursuit à la page suivante.}

\page{97}
un op. 1-diff. (\uncertain{rel.} : $X/S$)
\[
D : \underline{\omega} \longrightarrow L,
\]
\struck{\uncertain{satisfaisant}} et la restriction de $\sigma$ à $\underline{\omega}\otimes\underline{\omega} = \underline{\omega}^2 \subset \mathrm{P}^1_{X/S}(\omega)$
\[
\sigma_0 : \underline{\omega}^2 \to L \qquad \text{i.e. } \sigma_0 \in \Gamma(L\otimes\omega^{-2})
\]
est ainsi défini par
\[
D(f\omega) - f\,D(\omega) = \sigma_0(\omega\cdot df)
\]
($\omega\in\Gamma\underline{\omega}$, $f\in\Gamma\underline{O}_X$). Les pts où la section $\sigma$ canonique sont ceux où $\sigma_0$ s'annule, i.e. les \emph{pts de dégénérescence} de l'op. diff. $D$. On notera qu'alors \uncertain{l'ordre} de \uncertain{dég.}, il est \uncertain{automatiquement} pas \uncertain{tel} que \uncertain{$\sigma$} est épi --- que ce soit épi en ce pt signifie aussi que $L$ est engendré comme \uncertain{module} par les $D(\omega)$ ($\omega$ section loc. de $\underline{\omega}$) \uncertain{si} \uncertain{on} \uncertain{peut} \uncertain{vrai}, donc (si $X$ courbe alg. sur corps de base $k$) il y a une \uncertain{\ill{}} \uncertain{\ill{}} \uncertain{finie} $L'\subset L$ \uncertain{maximale} pour laquelle \uncertain{la} \ill{}
\note{les lignes 9 à 11 de la page sont pâles et serrées : lecture \uncertain{fragmentaire}.}

\ldots\ \uncertain{se factorise} $D$ (la plus petite \ldots) par \uncertain{\ill{}}, et \uncertain{quelle} \uncertain{que} \uncertain{soit} laquelle la condition \uncertain{est} vérifiée\ldots

On \uncertain{voit} donc (\uncertain{via} \struck{les} telles équations \ldots\ \uncertain{associées}, \emph{sur l'ouvert des pts} \emph{\uncertain{non} de dégénérescence}), une connexion \struck{\uncertain{projective}} \uncertain{singulière} qui \uncertain{\ill{}} \uncertain{rend} \uncertain{horizontale} \ldots\ \struck{Question : \uncertain{si} $\exists$ connexion}
\note{le bas de la page est lu par fragments ; la dernière ligne est biffée.}

\page{98}
Pour que la section donnée $\sigma$ de $P$ soit horizontale pour une conn. projective \uncertain{admissible} \uncertain{et} (\uncertain{a} \ldots) que l'ouvert $U$ des pts où $D$ est non dég. soit dense (\uncertain{en} chaque fibre de $X$ sur $S$ et b) que la section qui vient de \uncertain{définir} \add{sur $U$} \struck{\ill{}} --- la section \uncertain{\ill{}} du torseur $\underline{\mathrm{Conn}}_{\mathrm{proj}}$ sous $\underline{\omega}^2$ --- \struck{\ill{}} « soit \uncertain{régulière} sur tt $X$ » (i.e. se prolonge \add{en section régulière} sur tt $X$ (laquelle sera unique)). Cette condition est une condition purement locale en les pts de \struck{\uncertain{régularité}, et signifie} \add{\uncertain{\ill{}}} dégénérescence, et signifie qu'en tant que \struck{\ill{}} scindage de $\mathrm{P}^1_{X/S}(\underline{\omega})$, $\sigma$ \struck{\ill{}} doit avoir une singularité polaire au plus d'ordre 1, et la partie \struck{polaire} \struck{et la} principale polaire à l'ordre 1 près (\uncertain{\ill{}}) \uncertain{\ill{}} \uncertain{sections} nulles au pt \uncertain{envisagé} doit être déterminée (\uncertain{écrire} que $\lambda^2 + \frac{d\lambda}{dt}$ n'a pas de pôle en un pt où $\lambda$ \struck{\ill{}} \add{en a un} !). Donc, \uncertain{le} \uncertain{long} de la section canonique $\sigma$ de $P$, il y a un champ d'éléments de section d'ordre 2 transverses à $\sigma$, et les sections $\sigma'$ de $P$ \add{\uncertain{sur} $X$} \uncertain{lesquelles} il y a une connexion projective de $X$ rendant $\sigma'$ horizontale (laquelle sera unique) sont celles qui, en un pt \add{commun} d'intersection avec $\sigma$, ont \struck{un} pt d'ordre 2 celui du champ. [\emph{NB} des éléments de contact \uncertain{con.} d'ordre 1 \uncertain{à} \uncertain{traverser} \uncertain{de} \uncertain{façon} \uncertain{précise} qu'on prend un $\Omega$ inversible, \ldots\ la section de $\mathbb{P}(\mathrm{P}^1_{X/S}(\Omega))$ définie par la \ldots
\note{la fin de la page se poursuit à la page suivante ; plusieurs mots courts du passage ne sont lus que \uncertain{par conjecture}.}

\page{99}
structure d'extension de $\mathrm{P}^1_{X/S}(\Omega)$, \ldots\ --- plus gén., chaque fois qu'on a une ext. de $\Omega$ inv. par $\Omega\otimes\underline{\omega}$, $\underline{\omega} = \Omega^1_{X/S}$ inversible aussi\ldots]

Si l'op. diff.
\[
\omega \longrightarrow \underbrace{L \simeq \underline{O}_X}_{\text{\uncertain{trivialisé}}}
\]
est
\[
D(f\,dt) = a\frac{df}{dt} + b
\]
la \uncertain{condition} \uncertain{invariante} principale est $\frac{b}{a} = \lambda_t$ (ne dépend que du choix de l'uniformisante $t$) et la condition, en un pôle de $\lambda$ i.e. zéro de $a$ (qui \uncertain{donne} \uncertain{au} \uncertain{droit} \uncertain{par} \uncertain{des} zéros de $b$, \uncertain{car} $a$ et $b$ ne \uncertain{divisent} pas \uncertain{tous} \uncertain{simultanément}) est
\begin{itemize}
\item[a)] $\lambda_t$ n'a que \uncertain{pôles} simples
\item[b)] \struck{\ill{}} $\lambda_t$ a comme développement formel en $t$
\[
\lambda_t = \frac{2}{t}\,(+\ \text{pas de termes constants}) + \text{termes d'ordre} \geq 1
\]
\end{itemize}
\note{la parenthèse « $+$ pas de termes constants » est une lecture \uncertain{probable} ; le reste de la page est blanc.}

\page{100}
\note{feuille de brouillon, numérotée par lui « 1 » (cerclé) en haut à droite ; la note du coin supérieur gauche est biffée de traits obliques.}

\struck{$\mathcal{M}^{\uncertain{q}}_g \to \mathcal{M}_g$, $\mathcal{M}_g \subset$ \uncertain{bien} \uncertain{connu} \ldots\ (\uncertain{net à} \ill{} \uncertain{finies})}

\[
P = \mathbb{P}(\mathrm{P}^1_{X/S}(\underline{\omega})),\ \sigma,\ c
\longmapsto \sigma,\ c + \varpi_2 \qquad (\varpi_2 \neq 0)
\]
\uncertain{automorphisme} de $P$ transformant $c$ en $c+\varpi_2$

\uncertain{cas} \uncertain{Unique} \ldots\ \uncertain{automorphismes} ? Quels sont les \uncertain{automorphismes} $u$ de $P$ tels que $u(\sigma)$ \uncertain{\ill{}} de la forme $c+\varpi_2$ ??

$\ldots\ \omega \subset \mathcal{L} \subset \mathfrak{g}\ldots$

\emph{Sections de $P$} disjointes de $\sigma$ $=$ splittings de $\mathcal{L}$
\[
\begin{aligned}
&= \text{splittings de } \mathrm{P}^1_{X/S}(\underline{\omega})\\
&= \text{connexions sur } \omega
\end{aligned}
\]
\emph{ou} op. diff. \add{d'ordre 1} \struck{sur $\omega$, \uncertain{forme} \ldots} \add{sur $\omega$} à valeurs dans $\omega\otimes\omega = \omega^2$ \struck{Connexions, \uncertain{projectives}, s'interprétant comme} satisfaisant
\[
D(f\omega) - f\,D(\omega) = \omega\otimes df
\]
\[
D(f\,dt) = f\,\underbrace{D(dt)}_{\lambda} + dt\,df
\]
sections de $P$ non disjointes \struck{de $\sigma$} $=$ connexions de $\omega$ avec \struck{points} \emph{singularités} polaires.

Connexions projectives de $P$ $=$ \emph{\uncertain{conditions}} \uncertain{différentielles} du premier ordre sur sections de $P$ i.e. sur connexions de $\omega$, de \uncertain{la} forme \ldots ?

\[
0 \to \underbrace{\omega}_{} \subset \mathcal{L} \subset \mathfrak{g} \to 0
\qquad
\begin{aligned}
&v_s : \omega^{-1} \to \mathfrak{g}\\
&v_s \in \Gamma\,\mathfrak{g}\otimes\underline{\omega} = \underline{\omega}^2\oplus\underline{\omega}\oplus\underline{\omega}^0
\end{aligned}
\]
\note{à droite, d'une écriture très fine, « $\deg = 0+0+1$ » \uncertain{(?)}.}
\[
\underline{\omega} + \underline{O}_X\cdot s + v_s(\underline{\omega}^{-1}) \qquad c(s) \subset s\otimes\omega + v_s\cdot\underline{O}_X\ ?
\]
\[
\omega_a + \underline{O}_s + \omega' \qquad
\text{\struck{$\lambda$}}\ c(s) = \varphi_2 + s\otimes\varphi_1 + v_s\cdot\varphi_0
\quad \text{avec } \varpi_0 = 0\ ?
\]
\[
s = s_0 + a_1 \quad (a_1\in\Gamma\underline{\omega})
\]
\[
\begin{aligned}
(\exp(\Theta_{a_1})\otimes\mathrm{id}_{\omega})(c(s)) &= \varphi_2 + (\exp\Theta_{a_1} s)\otimes\varphi_1 + (\exp\Theta_{a_1} v_s)\,\varphi_0\\
&= \varphi_2 + s_0\otimes\varphi_1 + \varphi_0
\end{aligned}
\]
\note{les dernières lignes sont surchargées ; les indices de $\Theta$ et le dernier facteur de la troisième ligne (« $(\varphi_0 / s)$ ») sont \uncertain{incertains}. L'argument se poursuit au-delà de cette page.}

\end{document}
